\chapter{RESULTS AND DISCUSSION}

\ChapterSection{YOLO Training Results}

\ChapterSubsection{Training Dynamics}

The YOLO detector was trained for 100 epochs. The final recorded training time was 25,646.7
seconds, equivalent to 7.12 hours. The training run reached its best recorded epoch at epoch 100,
which indicates that the selected run did not terminate early before completing the configured
training schedule.

Figure~\ref{fig:yolo-training-dynamics} presents the training dynamics. The curve shows the
relationship between loss reduction and detection metrics across the training process. The model
improves substantially from the early epochs and stabilizes toward the later part of training.
The final epochs are important because the best mAP@50-95 value is recorded at the end of the
run, not in an earlier intermediate checkpoint.

\begin{figure}[H]
\centering
\includegraphics[width=\textwidth]{data/krishi_vaidya_yolo_data/outputs/figures/figure1_training_dynamics.pdf}
\caption{Training dynamics of the Krishi Vaidya YOLO detector}
\label{fig:yolo-training-dynamics}
\end{figure}

\begin{table}[H]
  \centering
  \caption{Summary of training dynamics and best-epoch performance for the YOLO detector}
  \label{tab:yolo-training-dynamics-summary}
  \begin{tabular}{lr}
    \toprule
    \textbf{Metric / Parameter} & \textbf{Value} \\
    \midrule
    Total epochs trained & $100$ \\
    Best training epoch & $100$ \\
    Input resolution (pixels) & $640 \times 640$ \\
    Total training time & $7.12\,\text{hours}$ \\
    \midrule
    \multicolumn{2}{l}{\textit{Loss residuals at best epoch}} \\
    Train box loss & $1.0627$ \\
    Train class loss & $0.1871$ \\
    Train DFL loss & $1.0493$ \\
    Validation box loss & $0.8855$ \\
    Validation class loss & $0.2681$ \\
    Validation DFL loss & $1.2635$ \\
    \midrule
    \multicolumn{2}{l}{\textit{Detection accuracy at best epoch}} \\
    Precision & $0.7172$ \\
    Recall & $0.7487$ \\
    F1-score & $0.7326$ \\
    mAP@50 & $0.7616$ \\
    mAP@50-95 & $0.6019$ \\
    \bottomrule
  \end{tabular}
\end{table}

\ChapterSubsection{Best-Epoch Performance}

Table~\ref{tab:yolo-training-dynamics-summary} summarizes the final and best-epoch values. The
model achieved precision of 0.7172, recall of 0.7487, F1-score of 0.7326, mAP@50 of 0.7616, and
mAP@50-95 of 0.6019. The recall is slightly higher than precision, which means the detector is
more inclined toward finding annotated objects than toward filtering out every uncertain
prediction. For the Krishi Vaidya workflow, this is a reasonable property for a first-stage visual
detector because missed plant regions can prevent later analysis, while uncertain detections can
still be filtered by thresholding, downstream validation, or user interaction.

The difference between mAP@50 and mAP@50-95 also matters. mAP@50 measures whether detected
objects overlap the ground truth at a moderate threshold, while mAP@50-95 requires stronger
localization across stricter IoU thresholds. The lower mAP@50-95 value indicates that localization
quality remains more difficult than basic object discovery, especially in a dataset with strongly
uneven class frequency and object scales.

\ChapterSection{YOLO Deployment and Threshold Results}

\ChapterSubsection{Threshold Selection Behavior}

The confidence-threshold sweep shows the expected precision-recall tradeoff. At the lowest tested
threshold of 0.05, recall reaches 0.7848, but precision is lower at 0.6748. At moderate thresholds,
precision rises while recall falls. The highest F1-score in the sweep is 0.7362 at a confidence
threshold of 0.10. This threshold gives precision of 0.7463 and recall of 0.7264.

\begin{figure}[H]
\centering
\includegraphics[width=\textwidth]{data/krishi_vaidya_yolo_data/outputs/figures/figure2_threshold_sweep.pdf}
\caption{Confidence-threshold sensitivity of the YOLO detector}
\label{fig:yolo-threshold-sweep}
\end{figure}

\begin{table}[H]
  \centering
  \caption{Sensitivity sweep of confidence thresholds on the YOLO baseline model}
  \label{tab:yolo-threshold-sensitivity-sweep}
  \small
  \resizebox{\textwidth}{!}{%
  \begin{tabular}{ccccccc}
    \toprule
    \textbf{Confidence threshold} & \textbf{Precision} & \textbf{Recall} & \textbf{F1-score} & \textbf{mAP@50} & \textbf{mAP@50-95} & \textbf{Fitness} \\
    \midrule
    0.05 & 0.6748 & 0.7848 & 0.7256 & 0.7225 & 0.5799 & 0.5799 \\
    \textbf{0.10} & \textbf{0.7463} & \textbf{0.7264} & \textbf{0.7362} & \textbf{0.6879} & \textbf{0.5613} & \textbf{0.5613} \\
    0.15 & 0.7710 & 0.6825 & 0.7241 & 0.6573 & 0.5441 & 0.5441 \\
    0.20 & 0.7923 & 0.6559 & 0.7177 & 0.6352 & 0.5314 & 0.5314 \\
    0.25 & 0.7729 & 0.6390 & 0.6996 & 0.6219 & 0.5227 & 0.5227 \\
    0.30 & 0.8016 & 0.6221 & 0.7006 & 0.6082 & 0.5138 & 0.5138 \\
    0.35 & 0.8084 & 0.6113 & 0.6962 & 0.5983 & 0.5065 & 0.5065 \\
    0.40 & 0.8241 & 0.6003 & 0.6946 & 0.5885 & 0.4986 & 0.4986 \\
    0.45 & 0.8348 & 0.5826 & 0.6863 & 0.5730 & 0.4860 & 0.4860 \\
    0.50 & 0.7480 & 0.5591 & 0.6399 & 0.5510 & 0.4712 & 0.4712 \\
    0.55 & 0.7579 & 0.5426 & 0.6324 & 0.5344 & 0.4587 & 0.4587 \\
    0.60 & 0.7725 & 0.5220 & 0.6230 & 0.5164 & 0.4466 & 0.4466 \\
    0.65 & 0.7978 & 0.4834 & 0.6020 & 0.4790 & 0.4221 & 0.4221 \\
    0.70 & 0.7997 & 0.4448 & 0.5717 & 0.4412 & 0.3962 & 0.3962 \\
    0.75 & 0.6166 & 0.4140 & 0.4954 & 0.4109 & 0.3728 & 0.3728 \\
    0.80 & 0.6143 & 0.3722 & 0.4635 & 0.3712 & 0.3412 & 0.3412 \\
    0.85 & 0.5433 & 0.3274 & 0.4086 & 0.3263 & 0.3044 & 0.3044 \\
    0.90 & 0.5429 & 0.2835 & 0.3725 & 0.2818 & 0.2643 & 0.2643 \\
    0.95 & 0.2727 & 0.0776 & 0.1208 & 0.0777 & 0.0729 & 0.0729 \\
    \bottomrule
  \end{tabular}%
  }
\end{table}

The threshold sweep also shows why a high confidence threshold is unsuitable for this detector.
At 0.95, recall falls to 0.0776 and F1-score falls to 0.1208. This would remove most detections
and make the detector ineffective as a first-stage model. Conversely, the 0.05 threshold has
higher recall but lower precision than the selected threshold. The 0.10 operating point therefore
provides a practical balance between keeping enough detections and controlling false positives.

\ChapterSubsection{Quantization Tradeoffs}

The deployment comparison evaluates the baseline PyTorch model and three TensorFlow Lite variants.
The baseline PyTorch model has the lowest latency in the recorded environment, with 12.18 ms per
image and 82.1 FPS. The TFLite variants are slower in this benchmark, but they represent the model
formats needed for lightweight deployment workflows.

\begin{figure}[H]
\centering
\includegraphics[width=\textwidth]{data/krishi_vaidya_yolo_data/outputs/figures/figure3_quantization_comparison.pdf}
\caption{Deployment comparison across baseline and quantized YOLO variants}
\label{fig:yolo-quantization-comparison}
\end{figure}

\begin{table}[H]
  \centering
  \caption{Deployment efficiency comparison across YOLO model variants}
  \label{tab:yolo-deployment-efficiency}
  \small
  \resizebox{\textwidth}{!}{%
  \begin{tabular}{lcccccccc}
    \toprule
    \textbf{Variant} & \textbf{Size (MB)} & \textbf{Compression} & \textbf{Device} & \textbf{Latency (ms)} & \textbf{FPS} & \textbf{Precision} & \textbf{Recall} & \textbf{F1} \\
    \midrule
    Baseline (PyTorch) & 5.97 & 1.00$\times$ & cuda:0 & 12.18 & 82.1 & 0.6929 & 0.7952 & 0.7402 \\
    FP32 (TFLite) & 11.72 & 0.51$\times$ & cuda:0 & 153.59 & 6.5 & 0.6930 & 0.7957 & 0.7402 \\
    FP16 (TFLite) & 5.90 & 1.01$\times$ & cuda:0 & 149.94 & 6.7 & 0.6903 & 0.7926 & 0.7377 \\
    INT8 (TFLite) & 3.20 & 1.87$\times$ & cuda:0 & 174.01 & 5.8 & 0.6774 & 0.7811 & 0.7268 \\
    \bottomrule
  \end{tabular}%
  }
\end{table}

The INT8 model is the smallest variant at 3.20 MB and gives the strongest compression ratio,
1.87$\times$ relative to the baseline size. This reduction has a measurable accuracy cost:
precision decreases from 0.6929 to 0.6774, recall decreases from 0.7952 to 0.7811, and F1-score
decreases from 0.7402 to 0.7268. FP16 preserves performance more closely, with F1-score of 0.7377
and a model size of 5.90 MB. The FP32 TFLite model matches the baseline F1-score but is larger
than the baseline model in this export.

These results show that quantization should not be selected only by size. For a mobile-first
system, INT8 is attractive because of its smaller footprint, but FP16 may be preferable when a
small accuracy loss is more important than maximum compression. Final selection should also
consider actual Android-device profiling, because the current latency values were measured on
\texttt{cuda:0}.

\ChapterSubsection{Per-Class Behavior}

Aggregate deployment metrics hide important class-level differences. Table~\ref{tab:yolo-per-class-f1-comparison}
shows that the best variant changes by class. The baseline PyTorch model gives the best F1-score
for rice leaf, rice panicle, tomato leaf, brassica leaf, and cabbage head. FP32 performs best for
rice grain cluster, maize ear, and tomato fruit. FP16 performs best for potato leaf. INT8 performs
best for maize leaf and cauliflower head.

\begin{table}[H]
  \centering
  \caption{Per-class F1-score comparison across YOLO model variants}
  \label{tab:yolo-per-class-f1-comparison}
  \small
  \resizebox{\textwidth}{!}{%
  \begin{tabular}{lcccc}
    \toprule
    \textbf{Class} & \textbf{Baseline (PyTorch)} & \textbf{FP32 (TFLite)} & \textbf{FP16 (TFLite)} & \textbf{INT8 (TFLite)} \\
    \midrule
    Rice leaf & \textbf{0.7750} & 0.7716 & 0.7207 & 0.7124 \\
    Rice panicle & \textbf{0.7799} & 0.7295 & 0.7277 & 0.7426 \\
    Rice grain cluster & 0.6911 & \textbf{0.7563} & 0.6980 & 0.7093 \\
    Maize leaf & 0.7139 & 0.7162 & 0.7137 & \textbf{0.7655} \\
    Maize ear & 0.6934 & \textbf{0.7736} & 0.7115 & 0.7363 \\
    Potato leaf & 0.7446 & 0.7219 & \textbf{0.7857} & 0.6977 \\
    Tomato leaf & \textbf{0.7793} & 0.7089 & 0.7336 & 0.6885 \\
    Tomato fruit & 0.7788 & \textbf{0.7852} & 0.7085 & 0.7083 \\
    Brassica leaf & \textbf{0.7876} & 0.7327 & 0.7240 & 0.7118 \\
    Cauliflower head & 0.7164 & 0.7058 & 0.7255 & \textbf{0.7361} \\
    Cabbage head & \textbf{0.7833} & 0.7252 & 0.7740 & 0.7245 \\
    \bottomrule
  \end{tabular}%
  }
\end{table}

This result suggests that quantization affects classes unevenly. The effect is not simply a
uniform decline from baseline to INT8. Some classes improve under a converted model, while others
lose performance. Because Krishi Vaidya targets multiple crop structures, deployment decisions
should therefore consider the classes most important to the intended use case, not only the
aggregate F1-score.

\ChapterSection{YOLO Error Analysis}

\ChapterSubsection{Confusion Patterns}

Figure~\ref{fig:yolo-confusion-matrix} presents the confusion matrix for the YOLO detector. The
confusion matrix is necessary because the dataset has severe class imbalance and because aggregate
metrics do not show whether failures are concentrated in a few classes.

\begin{figure}[H]
\centering
\includegraphics[width=0.92\textwidth]{data/krishi_vaidya_yolo_data/outputs/figures/figure4_confusion_matrix.pdf}
\caption{Confusion matrix of the YOLO detector}
\label{fig:yolo-confusion-matrix}
\end{figure}

\ChapterSubsection{Error Categories}

Table~\ref{tab:yolo-error-analysis} separates true positives, missed detections, and ghost
detections by class. The most significant failure is the rice panicle class, where 16,147 objects
are missed, corresponding to 92.9\% false negatives. Rice leaf also has a very high missed rate,
with 339 missed objects and only 7 true positives. Brassica leaf has 75 missed objects and no true
positives in the reported breakdown.

\begin{table}[H]
  \centering
  \caption{Error analysis breakdown by YOLO class}
  \label{tab:yolo-error-analysis}
  \small
  \resizebox{\textwidth}{!}{%
  \begin{tabular}{lrrr}
    \toprule
    \textbf{Class} & \textbf{True positives} & \textbf{False negatives (missed)} & \textbf{False positives (ghost)} \\
    \midrule
    \textbf{Rice panicle} & 1,238 & 16,147 (92.9\%) & 46 (3.6\%) \\
    Rice leaf & 7 & 339 (98.0\%) & 4 (36.4\%) \\
    Tomato fruit & 168 & 119 (41.5\%) & 9 (5.1\%) \\
    Rice grain cluster & 58 & 24 (29.3\%) & 68 (54.0\%) \\
    Brassica leaf & 0 & 75 (100.0\%) & 0 (0.0\%) \\
    Potato leaf & 129 & 19 (12.8\%) & 23 (15.1\%) \\
    Maize leaf & 990 & 10 (1.0\%) & 11 (1.1\%) \\
    Tomato leaf & 1,039 & 8 (0.8\%) & 4 (0.4\%) \\
    Maize ear & 3 & 1 (25.0\%) & 4 (57.1\%) \\
    Cabbage head & 45 & 0 (0.0\%) & 2 (4.3\%) \\
    Cauliflower head & 24 & 0 (0.0\%) & 0 (0.0\%) \\
    \bottomrule
  \end{tabular}%
  }
\end{table}

The error pattern shows that the detector is not equally reliable across classes. Maize leaf and
tomato leaf show strong behavior in this breakdown, with low missed-detection and ghost-detection
rates. Cabbage head and cauliflower head show no missed detections in the table, although their
absolute true-positive counts are small. In contrast, rice panicle and rice leaf have high
missed-detection rates, and rice grain cluster has a high ghost-detection rate. These results
indicate that class-specific improvements are more important than simply extending training time.

The rice panicle result deserves careful interpretation. The dataset contains 115,046 rice
panicle annotations, far more than any other class. Many panicle objects are likely small,
clustered, or visually dense, and the curation report also shows that many dropped objects were
related to out-of-bounds, truncated, or very small annotations. The high false-negative count
therefore indicates a concrete limitation of the current detector: it can learn useful aggregate
detection behavior, but dense or small structures remain difficult.

\ChapterSubsection{Summary of YOLO Findings}

The YOLO component provides a complete, evidence-backed object-detection subsystem for Krishi
Vaidya. The trained detector achieved F1-score of 0.7326, mAP@50 of 0.7616, and mAP@50-95 of
0.6019 at the best recorded epoch. The threshold sweep identifies 0.10 as the best F1 operating
point among the tested confidence thresholds. The deployment comparison shows that TFLite variants
preserve useful detection performance, with INT8 providing the smallest model size and FP16
providing a closer accuracy-size compromise.

The main limitation is class-specific reliability. The detector performs much better for some
classes than others, and dense or underrepresented structures require further improvement. For the
larger Krishi Vaidya system, this means the YOLO model is suitable as a visual detection component,
but its outputs should be treated as supporting evidence rather than as the sole basis for disease
diagnosis.

\ChapterSection{VLM Training and Resource Results}

\ChapterSubsection{Fine-Tuning Performance}

The VLM training run completed 3 epochs and 969 training steps using Qwen2.5-VL-3B-Instruct
as the base model. Table~\ref{tab:vlm-training-summary}
summarizes the training outcome. The final reported evaluation loss is 0.0146, and the mean token
accuracy is 0.9967. The trainer-reported train loss is 0.1048, while the last logged train loss is
0.0097. These values show that the supervised fine-tuning process fitted the training objective
strongly.

\begin{figure}[H]
\centering
\includegraphics[width=\textwidth]{data/krishi_vaidya_vlm_data/outputs/figures/figure1_training_dynamics.pdf}
\caption{Training dynamics of the Qwen2.5-VL LoRA fine-tuning run}
\label{fig:vlm-training-dynamics}
\end{figure}

\begin{table}[H]
  \centering
  \caption{Training summary for Qwen2.5-VL LoRA fine-tuning}
  \label{tab:vlm-training-summary}
  \begin{tabular}{lr}
    \toprule
    \textbf{Metric / setting} & \textbf{Value} \\
    \midrule
    Base model & Qwen/Qwen2.5-VL-3B-Instruct \\
    Epochs completed & 3.00 \\
    Train steps & 969 \\
    Last logged train loss & 0.0097 \\
    Trainer-reported train loss & 0.1048 \\
    Final eval loss & 0.0146 \\
    Mean token accuracy & 0.9967 \\
    Runtime & 2.58 h \\
    Samples/sec & 1.664 \\
    Learning rate & 2.00e-05 \\
    Effective batch & 4 x 4 \\
    LoRA rank / alpha & 16 / 32 \\
    Quantization & 4-bit nf4 \\
    \bottomrule
  \end{tabular}
\end{table}

The training result must be interpreted carefully. The low evaluation loss and high token accuracy
indicate that the model learned the supervised answer format on the validation objective. They do
not, by themselves, prove that the model will produce parseable and correct disease labels during
free generation on the test split. For that reason, the VLM evaluation includes a separate
prediction-metrics study, discussed in Section~\ref{sec:vlm-prediction-results}.

\ChapterSubsection{Runtime and Memory Behavior}

The resolved runtime used an NVIDIA L4 GPU with CUDA 12.8 and compute capability 8.9. The run
used \texttt{flash\_attention\_2}, Liger kernel support, bfloat16, and TF32. The memory telemetry
shows peak allocated GPU memory of 12.56 GB and peak reserved memory of 13.08 GB, within the
available 22.03 GB GPU memory.

\begin{figure}[H]
\centering
\includegraphics[width=\textwidth]{data/krishi_vaidya_vlm_data/outputs/figures/figure2_memory_runtime.pdf}
\caption{Runtime and memory behavior of the VLM training run}
\label{fig:vlm-memory-runtime}
\end{figure}

\begin{table}[H]
  \centering
  \caption{Runtime and resource envelope for VLM fine-tuning}
  \label{tab:vlm-runtime-resources}
  \begin{tabular}{lr}
    \toprule
    \textbf{Metric / setting} & \textbf{Value} \\
    \midrule
    GPU & NVIDIA L4 \\
    CUDA / capability & 12.8 / 8.9 \\
    GPU memory & 22.03 GB \\
    Peak allocated & 12.56 GB \\
    Peak reserved & 13.08 GB \\
    Peak RAM used & 40.6\% \\
    Attention backend & flash\_attention\_2 \\
    Liger kernel & True \\
    torch.compile & False \\
    bf16 / tf32 & True / True \\
    \bottomrule
  \end{tabular}
\end{table}

The resource profile supports the technical feasibility of adapter fine-tuning a 3B-parameter VLM
under a single-GPU setup. The use of 4-bit NF4 quantization, LoRA adapters, bfloat16 computation,
and efficient attention was important for keeping the run inside the available GPU memory.

\ChapterSection{VLM Prediction Results}
\label{sec:vlm-prediction-results}

\ChapterSubsection{Dataset Split and Class Distribution}

The validation study confirms that the prepared Qwen export contains 6,647 records and 13,294
image references. The train split contains 5,147 Qwen records, the validation split contains 702
records, and the test split contains 798 records. The image count is exactly twice the record
count in each split, as reported by the validation output.

\begin{figure}[H]
\centering
\includegraphics[width=\textwidth]{data/krishi_vaidya_vlm_data/outputs/figures/figure3_dataset_composition.pdf}
\caption{Dataset composition for the VLM Qwen export}
\label{fig:vlm-dataset-composition}
\end{figure}

\begin{table}[H]
  \centering
  \caption{Dataset split validation summary for VLM records}
  \label{tab:vlm-dataset-splits}
  \begin{tabular}{lrr}
    \toprule
    \textbf{Split} & \textbf{Qwen records} & \textbf{Images found} \\
    \midrule
    Train & 5,147 & 10,294 \\
    Validation & 702 & 1,404 \\
    Test & 798 & 1,596 \\
    Total & 6,647 & 13,294 \\
    \bottomrule
  \end{tabular}
\end{table}

\begin{table}[H]
  \centering
  \caption{Qwen record distribution by VLM class}
  \label{tab:vlm-class-distribution}
  \begin{tabular}{lr}
    \toprule
    \textbf{Class} & \textbf{Qwen records} \\
    \midrule
    Maize: Common Rust & 412 \\
    Maize: Healthy & 800 \\
    Maize: Northern Leaf Blight & 800 \\
    Potato: Early Blight & 800 \\
    Potato: Healthy & 72 \\
    Potato: Late Blight & 490 \\
    Tomato: Early Blight & 477 \\
    Tomato: Healthy & 800 \\
    Tomato: Late Blight & 800 \\
    Tomato: Leaf Mold & 396 \\
    Tomato: Septoria Leaf Spot & 800 \\
    \bottomrule
  \end{tabular}
\end{table}

The class distribution is partially capped but still uneven. Several classes contain 800 records,
while potato healthy contains only 72 records. This imbalance affects interpretation of macro
metrics and per-class recall. The test split is also uneven: maize common rust alone contributes
412 of the 798 test examples, which makes the model's failure on that class especially important
for the aggregate accuracy.

\ChapterSubsection{Prediction Summary}

The test-set prediction metrics reveal a clear gap between supervised fine-tuning loss and usable
structured prediction. Table~\ref{tab:vlm-prediction-summary} reports accuracy of 0.1303, macro
precision of 0.5232, macro recall of 0.1834, and macro F1-score of 0.2465 over 798 test examples.
The most important result is that 540 predictions, or 67.7\%, are parse-error predictions.

\begin{figure}[H]
\centering
\includegraphics[width=\textwidth]{data/krishi_vaidya_vlm_data/outputs/figures/figure4_prediction_performance.pdf}
\caption{Prediction performance of the VLM on the test split}
\label{fig:vlm-prediction-performance}
\end{figure}

\begin{table}[H]
  \centering
  \caption{Prediction summary on the VLM test split}
  \label{tab:vlm-prediction-summary}
  \begin{tabular}{lr}
    \toprule
    \textbf{Metric} & \textbf{Value} \\
    \midrule
    Split & test \\
    Examples & 798 \\
    Accuracy & 0.1303 \\
    Macro precision & 0.5232 \\
    Macro recall & 0.1834 \\
    Macro F1 & 0.2465 \\
    Parse-error predictions & 540 (67.7\%) \\
    \bottomrule
  \end{tabular}
\end{table}

This result is a central limitation of the current VLM implementation. The model was not judged
only by whether the generated text looked plausible; it was judged by whether the generated output
could be parsed into a canonical label for downstream use. A model response that cannot be parsed
into \texttt{canonical\_label} cannot reliably drive the backend diagnosis or mobile-app result
flow. Therefore, the parse-error rate directly reduces the practical diagnostic usability of the
current VLM.

\ChapterSubsection{Per-Class Prediction Metrics}

The per-class metrics in Table~\ref{tab:vlm-per-class-metrics} show that the VLM behaves unevenly
across classes. Maize healthy obtains the highest F1-score at 0.5556. Tomato leaf mold reaches
0.5143, potato late blight reaches 0.4681, tomato septoria leaf spot reaches 0.4138, and tomato
late blight reaches 0.3939. Several classes perform poorly: maize common rust has 412 false
negatives and no true positives, potato healthy has no true positives, and tomato early blight has
only one true positive.

\begin{table}[H]
  \centering
  \caption{Per-class prediction metrics for the VLM test split}
  \label{tab:vlm-per-class-metrics}
  \scriptsize
  \resizebox{\textwidth}{!}{%
  \begin{tabular}{lrrrrrrr}
    \toprule
    \textbf{Class} & \textbf{Support} & \textbf{Precision} & \textbf{Recall} & \textbf{F1} & \textbf{TP} & \textbf{FN} & \textbf{FP} \\
    \midrule
    Maize: Common Rust & 412 & 0.0000 & 0.0000 & 0.0000 & 0 & 412 & 0 \\
    Maize: Healthy & 50 & 0.9091 & 0.4000 & 0.5556 & 20 & 30 & 2 \\
    Maize: Northern Leaf Blight & 55 & 0.1296 & 0.3818 & 0.1935 & 21 & 34 & 141 \\
    Potato: Early Blight & 45 & 1.0000 & 0.1778 & 0.3019 & 8 & 37 & 0 \\
    Potato: Healthy & 5 & 0.0000 & 0.0000 & 0.0000 & 0 & 5 & 0 \\
    Potato: Late Blight & 35 & 0.9167 & 0.3143 & 0.4681 & 11 & 24 & 1 \\
    Tomato: Early Blight & 26 & 0.3333 & 0.0385 & 0.0690 & 1 & 25 & 2 \\
    Tomato: Healthy & 52 & 1.0000 & 0.1731 & 0.2951 & 9 & 43 & 0 \\
    Tomato: Late Blight & 50 & 0.8125 & 0.2600 & 0.3939 & 13 & 37 & 3 \\
    Tomato: Leaf Mold & 25 & 0.9000 & 0.3600 & 0.5143 & 9 & 16 & 1 \\
    Tomato: Septoria Leaf Spot & 43 & 0.8000 & 0.2791 & 0.4138 & 12 & 31 & 3 \\
    Parse Error & 0 & 0.0000 & 0.0000 & 0.0000 & 0 & 0 & 540 \\
    Maize: Early Leaf Spot & 0 & 0.0000 & 0.0000 & 0.0000 & 0 & 0 & 1 \\
    \bottomrule
  \end{tabular}%
  }
\end{table}

The high false-positive count for maize northern leaf blight is also significant. The model
predicts maize northern leaf blight 162 times, but only 21 of those predictions are correct. This
creates 141 false positives for that label. The model also emits one prediction for maize early
leaf spot, which is not a supported class in the reported test distribution. This indicates that
the VLM can generate labels outside the expected class set unless output constraints are enforced.

\ChapterSubsection{Confusion Matrix and Prediction Distribution}

The confusion matrix in Figure~\ref{fig:vlm-confusion-matrix} visualizes how predictions are
distributed across expected labels. The large parse-error column dominates the evaluation, which
is consistent with the prediction summary. This confirms that the main failure mode is not merely
confusion between two similar disease classes; it is failure to produce a parseable canonical
answer in many cases.

\begin{figure}[H]
\centering
\includegraphics[width=0.95\textwidth]{data/krishi_vaidya_vlm_data/outputs/figures/figure5_confusion_matrix.pdf}
\caption{Confusion matrix for VLM prediction results}
\label{fig:vlm-confusion-matrix}
\end{figure}

\begin{figure}[H]
\centering
\includegraphics[width=\textwidth]{data/krishi_vaidya_vlm_data/outputs/figures/figure6_inference_diagnostics.pdf}
\caption{Inference diagnostics for VLM predictions}
\label{fig:vlm-inference-diagnostics}
\end{figure}

\begin{table}[H]
  \centering
  \caption{Prediction distribution and median latency for the VLM test split}
  \label{tab:vlm-prediction-distribution}
  \small
  \resizebox{\textwidth}{!}{%
  \begin{tabular}{lrrr}
    \toprule
    \textbf{Predicted label} & \textbf{Count} & \textbf{Accuracy} & \textbf{Median latency} \\
    \midrule
    Parse Error & 540 & 0.0000 & 7.32 s \\
    Maize: Northern Leaf Blight & 162 & 0.1296 & 6.67 s \\
    Maize: Healthy & 22 & 0.9091 & 5.32 s \\
    Tomato: Late Blight & 16 & 0.8125 & 5.92 s \\
    Tomato: Septoria Leaf Spot & 15 & 0.8000 & 6.50 s \\
    Potato: Late Blight & 12 & 0.9167 & 6.15 s \\
    Tomato: Leaf Mold & 10 & 0.9000 & 5.71 s \\
    Tomato: Healthy & 9 & 1.0000 & 5.32 s \\
    Potato: Early Blight & 8 & 1.0000 & 6.23 s \\
    Tomato: Early Blight & 3 & 0.3333 & 5.90 s \\
    Maize: Early Leaf Spot & 1 & 0.0000 & 5.88 s \\
    \bottomrule
  \end{tabular}%
  }
\end{table}

The prediction distribution shows that parse errors are not only frequent but also slower than
most successful label categories, with a median latency of 7.32 seconds. Non-parse-error labels
generally have median latencies between 5.32 and 6.67 seconds. This suggests that the current
generation path is expensive enough to require careful backend integration and that output
formatting failures are costly in both accuracy and time.

\ChapterSubsection{Summary of VLM Findings}

The VLM work demonstrates that Krishi Vaidya includes a real fine-tuned vision-language component,
not only a detector or mobile interface. The model was fine-tuned with LoRA adapters on
Qwen2.5-VL-3B-Instruct using 4-bit NF4 quantization and completed the training run within the
available NVIDIA L4 resource envelope. The training objective was fitted successfully, as shown by
the low final evaluation loss and high token accuracy.

However, the prediction results show that the current VLM is not yet reliable as an autonomous
structured disease classifier. The major limitation is parseability: 540 of 798 test predictions
could not be parsed into a canonical label. The model also shows class-level weaknesses, especially
for maize common rust, potato healthy, and tomato early blight, and it can emit labels outside the
expected test label set. For the larger Krishi Vaidya system, the VLM should therefore be described
as a fine-tuned diagnostic component with clear current limitations. Its next improvement step is
not simply more training; it requires stricter output control, improved prompt consistency,
possibly constrained decoding or post-processing, and targeted correction of weak classes.

\ChapterSection{Backend Implementation Results}

\ChapterSubsection{Implemented API Surface}

The backend implementation provides the server-side application layer required by Krishi Vaidya.
The verified codebase contains eight main API domains: diagnosis, authentication, users, scans,
crops, tasks, notes, and advisories. These domains cover the system's core workflow: a farmer can
authenticate, maintain profile information, upload crop images, preserve scan history, manage crop
plots, attach notes, track tasks, and receive or manage advisories. The backend also exposes a
health endpoint and development-time Swagger documentation.

\begin{table}[H]
\centering
\caption{Backend implementation coverage by domain}
\label{tab:backend-implementation-coverage}
\small
\resizebox{\textwidth}{!}{%
\begin{tabular}{@{}llll@{}}
\toprule
\textbf{Domain} & \textbf{HTTP coverage} & \textbf{Persistence} & \textbf{Important behavior} \\
\midrule
Diagnosis & POST & External VLM and Google AI calls & Validates image, calls VLM, generates advisory JSON \\
Authentication & POST & User, OTP, refresh-token models & Register, login, refresh, logout, OTP, password reset \\
Users & GET, PATCH, PUT, DELETE, POST & User model & Profile, password, profile picture, push token management \\
Scans & GET, POST, PATCH, DELETE & Scan and scan-resolution models & Upload, filter, group, resolve, follow-up, status tracking \\
Crops & GET, POST, PATCH, DELETE & Crop-plot model & Plot lifecycle, archive state, stage update, health update \\
Tasks & GET, POST, PATCH, DELETE & Crop-task model & Pending, overdue, upcoming, completion, batch creation \\
Notes & GET, POST, PATCH, DELETE & Crop-note model & Plot notes, scan-linked notes, image attachments \\
Advisories & GET, POST, PATCH, DELETE & Advisory model & Active advisories, dismiss, undismiss, batch creation \\
\bottomrule
\end{tabular}%
}
\end{table}

The implemented API surface is broader than a diagnosis-only prototype. It supports both disease
analysis and crop-management records. This matters for the final system because the mobile
application needs more than a single prediction endpoint: it must keep user-specific farm context,
historical scans, and advisory state. The backend provides this context through authenticated,
user-scoped resources rather than through anonymous image processing alone.

\ChapterSubsection{Layered Implementation Outcome}

The final backend structure follows a consistent request path. Incoming requests enter the Express
application, pass through middleware, reach a route group, execute controller validation, call a
service method, and then use repository methods for persistence when required. The same pattern is
used across scans, crops, tasks, notes, and advisories, which reduces implementation inconsistency
between domains.

\begin{table}[H]
\centering
\caption{Backend implementation layers and their realized responsibilities}
\label{tab:backend-layers-results}
\small
\resizebox{\textwidth}{!}{%
\begin{tabular}{@{}lll@{}}
\toprule
\textbf{Layer} & \textbf{Implemented files} & \textbf{Realized responsibility} \\
\midrule
Configuration & \texttt{src/config} & Environment validation, MongoDB, Redis, Cloudinary, Swagger, runtime options \\
Application shell & \texttt{src/app.ts} & Middleware registration, health endpoint, route mounting, service initialization \\
Container & \texttt{src/container.ts} & Explicit dependency construction for repositories, services, controllers, routes \\
Middleware & \texttt{src/middlewares} & Authentication, CORS, Helmet, rate limits, upload parsing, validation, logging, errors \\
Routes & \texttt{src/routes} & Versioned API grouping and Swagger-annotated endpoint definitions \\
Controllers & \texttt{src/controllers} & Request validation, parameter extraction, response formatting \\
Services & \texttt{src/services} & Domain rules, image upload orchestration, diagnosis pipeline, advisory generation \\
Repositories & \texttt{src/repositories} & Mongoose queries, pagination, ownership filters, soft-delete access patterns \\
Models & \texttt{src/models} & MongoDB schemas, indexes, validation rules, document methods \\
\bottomrule
\end{tabular}%
}
\end{table}

This separation is useful for thesis interpretation because it shows that the backend was designed
as a maintainable application component, not as a collection of unrelated Express handlers. The
service and repository layers also make testing more practical because domain behavior can be
tested below the HTTP layer.

\ChapterSection{Backend Domain Discussion}

\ChapterSubsection{Authentication and User Management}

The authentication implementation supports registration, login, access-token refresh, logout,
logout from all devices, OTP sending, OTP verification, forgot-password, and password reset. The
password policy requires at least eight characters and at least one uppercase letter, one lowercase
letter, and one number. Refresh tokens are not stored as plain tokens; they are represented through
hashed records with lookup identifiers and revocation state. This is an appropriate design choice
because refresh tokens are long-lived compared with access tokens and should not be recoverable
from the database as usable secrets.

User management includes profile retrieval and update, profile-picture upload, password change,
account deletion, and push-token add/remove operations. The profile schema includes bilingual and
Nepal-specific fields such as Nepali name, preferred language, ecological region, province,
district, municipality, ward, latitude, and longitude. It also includes theme and color-blind mode
preferences, which shows that the backend stores application-level preference state rather than
only authentication data.

\ChapterSubsection{Scan Lifecycle}

The scan domain is the most complete persistence-backed workflow in the backend. A scan stores the
uploaded image URI, crop metadata, status, result object, notes, tags, offline flag, client-created
timestamp, deletion state, and optional processing timestamps. The scan service uploads images to
Cloudinary before creating scan records. It also supports scan retrieval, pagination, filtering by
status and crop type, soft deletion, batch deletion, grouped views, related scans, status checks,
scan resolution, resolution retrieval, follow-up linkage, pending-scan retrieval, offline-scan
retrieval, offline scan synchronization, result updates, and failed-scan marking.

The scan-resolution model adds an important post-diagnosis layer. A completed scan can be marked
as treated, resolved naturally, crop removed, or false positive. Treatment text, notes, and
follow-up scan identifiers can be stored. This makes the backend suitable for tracking what happens
after a disease prediction, which is important for an agricultural advisory system because diagnosis
alone does not measure whether the farmer acted on the recommendation or whether the crop condition
improved.

\ChapterSubsection{Crop, Task, Note, and Advisory Management}

Crop plots represent the farmer's managed crop areas. Each plot stores crop type, variety, region,
district, optional location, area value and unit, planting date, expected harvest date, growth
stage, stage progress, health status, archive state, and soft-delete state. The crop routes support
creation, retrieval, filtering, archive/unarchive actions, growth-stage updates, and health-status
updates.

Tasks provide a scheduling layer around crop management. The task model supports categories such as
watering, fertilizing, spraying, weeding, pruning, harvesting, soil preparation, scouting, and other
tasks. Each task has a due date, completion state, priority, and optional auto-generated flag. The
implemented routes expose pending, overdue, upcoming, plot-specific, batch, completion, and
uncompletion operations.

Notes provide a record-keeping layer. A note is linked to a crop plot and can optionally be linked
to a scan. It stores note type, title, body, image URIs, public image identifiers, tags, and
soft-delete state. The implementation supports adding and removing images after note creation,
which is important because farmers may add evidence over time rather than at the exact moment the
note is first created.

Advisories provide a notification-like layer for crop decisions. The advisory model supports
weather alerts, irrigation needs, fertilizer reminders, pest risk, disease risk, harvest windows,
spray windows, frost warnings, and stage transitions. Advisories include severity, bilingual title
and description fields, optional action labels, generated time, dismiss state, and expiry. The
model uses an expiry index, which allows expired advisories to be cleaned up automatically.

\ChapterSection{Backend Diagnosis Integration Discussion}

\ChapterSubsection{Structured Diagnosis Response}

The backend diagnosis pipeline produces a structured response rather than returning raw model text
to the mobile app. The response contains disease name, Nepali disease name, scientific name, crop,
confidence, confidence tier, out-of-distribution flag, symptom descriptions, advisory details,
processing time, request identifier, timestamp, and model version. The advisory object contains
treatment steps, prevention tips, safety warnings, organic alternatives, and source URLs.

This response structure is important because the mobile application needs stable fields for display
and decision support. It also reduces the risk of presenting unstructured model text directly to a
farmer. The backend still uses the VLM output, but it places a second structured-generation and
schema-validation step after the VLM to produce an application-facing advisory response.

\ChapterSubsection{Confidence and Out-of-Distribution Handling}

The backend converts confidence into four confidence tiers: high, medium, low, and
out-of-distribution. If the advisory generator marks a case as out of distribution, or if confidence is
below 0.30, the backend assigns the \texttt{ood} tier. Confidence values from generated outputs are
clamped between 0 and 1. Alternative diagnoses are included only when the confidence tier is medium
or low. This design avoids overloading high-confidence responses with unnecessary alternatives and
allows uncertain predictions to carry additional possibilities.

The design is sound, but it depends on the reliability of upstream VLM prediction and downstream
advisory JSON generation. The VLM prediction results in this manuscript show high parse-error
rates during independent evaluation. Therefore, the backend's structured advisory layer should be
understood as an integration safeguard, not as proof that the current VLM model is already
diagnostically reliable in all cases.

\ChapterSection{Backend Verification and Limitations}

\ChapterSubsection{Automated Test Coverage}

The available automated tests focus mainly on scan behavior. Unit tests cover scan creation,
offline-scan creation, owner-scoped retrieval, filtering, updating, soft deletion, statistics,
resolution creation, duplicate-resolution prevention, resolution retrieval, follow-up linkage, and
grouping by crop, disease, and month. Integration tests cover scan route authentication behavior,
paginated retrieval, filtering, statistics, grouping validation, lookup by identifier, update,
soft deletion, batch deletion, status retrieval, resolution creation, duplicate-resolution
protection, resolution retrieval, follow-up linkage, and related-scan lookup.

\begin{table}[H]
\centering
\caption{Implemented backend test focus}
\label{tab:backend-test-focus}
\small
\resizebox{\textwidth}{!}{%
\begin{tabular}{@{}lll@{}}
\toprule
\textbf{Test area} & \textbf{Evidence in repository} & \textbf{Covered behavior} \\
\midrule
Scan service unit tests & \texttt{tests/unit/services/scanService.test.ts} & Creation, ownership, filtering, update, delete, stats, resolution, follow-up, grouping \\
Scan route integration tests & \texttt{tests/integration/routes/scanRoutes.test.ts} & Auth requirement, pagination, filters, validation, status, resolution, related scans \\
Test database setup & \texttt{tests/setup.ts} & MongoDB memory server, database clearing, test environment variables \\
\bottomrule
\end{tabular}%
}
\end{table}

The test coverage is useful but not complete. It provides evidence that the scan workflow was
treated seriously, but it does not yet provide equivalent test evidence for every route group.
Authentication, crop, task, note, advisory, and diagnosis integration tests should be expanded
before claiming complete backend verification.

\ChapterSubsection{Current Backend Limitations}

Several limitations remain. First, the backend section currently has an architecture figure, but
it still lacks captured API screenshots, Swagger screenshots, request-response examples from a
running deployment, and deployment metrics. Second, the scan test suite is stronger than the test
coverage for the other domains. Third, Redis is initialized with BullMQ-compatible options, but the
visible implementation does not yet demonstrate a completed background-job pipeline for diagnosis
processing. Fourth, the diagnosis route uses external VLM and Google AI services; therefore, final
system reliability depends on those services being reachable, correctly configured, and returning
valid schemas. Fifth, direct production-load testing and mobile-to-backend latency measurements
have not yet been added to the manuscript evidence.

These limitations do not invalidate the backend implementation. They define what should be added
before the final manuscript claims production readiness. The current backend can be described as a
substantial implemented API layer with verified architecture, domain models, route groups,
security controls, image handling, and diagnosis integration, while final deployment evidence and
broader automated tests remain future documentation tasks.

\ChapterSection{Mobile App Implementation Results}

\ChapterSubsection{Implemented Application Surface}

The mobile application implementation shows that Krishi Vaidya has progressed beyond a model-only
prototype. The app contains a complete navigation shell, onboarding flow, scan workflow, crop
dashboard, history view, settings area, notification handling, export features, and local
persistence. The implemented source also contains development fallbacks for web execution and
native-specific paths for camera and TFLite functionality. This is important because the final
system must support actual farmer workflows rather than only demonstrate that an image can be sent
to a model.

\begin{table}[H]
\centering
\caption{Mobile application implementation coverage}
\label{tab:mobile-implementation-coverage}
\small
\resizebox{\textwidth}{!}{%
\begin{tabular}{@{}lll@{}}
\toprule
\textbf{Area} & \textbf{Implemented evidence} & \textbf{Result for the manuscript} \\
\midrule
Application shell & Expo config, root layout, providers, typed routes & Production-shaped app skeleton and controlled startup sequence \\
Onboarding & Welcome, language, permission, crop, demo screens & First-run user preparation before main workflows \\
Scan workflow & Capture, preview, editor, analysis, final result screens & End-to-end image-to-report user path in the app \\
Crop management & Plot dashboard, add/edit crop, tasks, notes, crop weather & Long-term agricultural record keeping beyond diagnosis \\
History and export & Scan history, grouping/filtering, export and result report support & Retention and reuse of diagnosis outputs \\
Settings and accessibility & Language, region, model precision, notifications, accessibility, profile & User personalization and operational control \\
Offline support & SQLite schema, offline queue, sync queue, network monitor, sync manager & Local-first behavior with queued work when connectivity is weak \\
\bottomrule
\end{tabular}%
}
\end{table}

The application surface is therefore consistent with the thesis objective of building an
integrated crop-disease support system. However, the evidence currently comes mainly from source
inspection and generated architecture documentation. Final manuscript completion should still add
screenshots from a running build, device-specific performance measurements, and user-flow captures
for the important screens.

\ChapterSubsection{Scan and Diagnosis Workflow Result}

The scan workflow is one of the strongest mobile implementation areas. The native capture view is
not limited to opening the camera. It includes selected-crop context, camera permission handling,
gallery selection, capture controls, model overlay state, and the ability to pass pre-detection
information forward. The preview stage checks image quality before analysis, and the analysis stage
uses a visible phase model. This gives the app a controlled interaction sequence: capture or select
an image, review image quality, optionally refine the image, run detection, submit for diagnosis,
and display a structured report.

The cloud diagnosis path is directly aligned with the backend diagnosis endpoint. The mobile API
client posts multipart image data to \texttt{/api/v1/diagnose} and accepts an optional crop-name
hint. The diagnosis hook then adapts behavior to network status. Online devices attempt cloud
diagnosis. Slow connections attempt cloud diagnosis with a five-second fallback threshold. Offline
devices queue the scan and attempt cached advisory output. This design is appropriate for
agricultural use in Nepal because field connectivity can be unstable, and the application should
not fail completely when the server is temporarily unreachable.

The result workflow is also broader than a plain prediction label. The result view model parses
diagnosis output for disease fields, advisory content, treatment timeline steps, organic
alternatives, voice output, feedback, reminder setup, share-card generation, and PDF export. This
matters because advisory delivery and follow-up actions are the parts of the system that make the
diagnosis usable to a farmer.

\ChapterSubsection{Local Data and Offline Result}

The local data layer provides a substantial offline foundation. Scan history, queued uploads, sync
operations, crop plots, crop tasks, crop notes, crop advisories, custom categories, user settings,
achievements, reminders, and notifications are represented in SQLite. The app also includes store
refresh logic that loads these records into UI-facing state during startup. This allows the app to
display prior records and user preferences even when network access is unavailable.

The sync implementation has two important outcomes. First, it provides concrete infrastructure for
offline-to-online recovery: pending queue items, batch processing, retry state, network monitoring,
app-resume triggers, and periodic synchronization. Second, it makes visible where final integration
work remains. The sync manager source marks some generic pending database changes as synced until
corresponding backend table-specific push handlers are implemented. Therefore, the correct
manuscript interpretation is that offline queueing and sync infrastructure are implemented, while
complete two-way synchronization for every local table remains a future integration task.

\ChapterSubsection{Crop Management and Farmer Record-Keeping Result}

The mobile app includes a crop-management subsystem that complements the diagnosis workflow. The
crop dashboard aggregates active plots, upcoming tasks, recent notes, active advisories, crop
statistics, network status, and weather-related context. Crop repositories support creating,
updating, archiving, filtering, and retrieving crop plots. Task and note repositories support
record keeping around farm activities. The data-catalog repository provides read-only crop
reference information from fixture files.

This is a strong design decision because a farmer-facing disease application should not treat each
scan as isolated. Crop condition, planting date, crop stage, task history, notes, weather context,
and advisory reminders all provide context for later decisions. The implementation therefore
supports a richer thesis claim: Krishi Vaidya is a mobile crop-health management system with an
AI-assisted diagnosis workflow, not only a crop-image classifier.

\ChapterSubsection{Model Deployment and Mobile Integration Result}

The mobile app includes TFLite model assets in int8, fp16, and fp32 precision variants and exposes
model precision as a user-facing setting. The TFLite service attempts GPU loading and falls back to
CPU loading, which is a practical deployment pattern for Android devices with different hardware
support. It also contains an inference benchmark method, so the app has a source-level path for
measuring load and inference time by precision.

\begin{table}[H]
\centering
\caption{Mobile model deployment evidence and integration notes}
\label{tab:mobile-model-deployment}
\small
\resizebox{\textwidth}{!}{%
\begin{tabular}{@{}lll@{}}
\toprule
\textbf{Item} & \textbf{Implemented evidence} & \textbf{Interpretation} \\
\midrule
TFLite assets & int8, fp16, and fp32 model files under mobile assets & Multiple mobile deployment variants are packaged \\
Runtime loader & Singleton TFLite service with selected precision & The app can choose a precision variant at runtime \\
Delegate strategy & Android GPU attempt with CPU fallback & Device heterogeneity is considered in the implementation \\
Benchmark hook & Load-time and inference-time return values in service method & Mobile profiling can be performed from existing code paths \\
Taxonomy alignment & Mobile labels differ from YOLO manuscript detector taxonomy & Final integration review is required before claiming full detector parity \\
\bottomrule
\end{tabular}%
}
\end{table}

The last row is important. The mobile source currently defines disease labels in the TFLite
service, while the YOLO section of this manuscript documents the trained detector from its own
training assets. This does not mean the mobile implementation is invalid, but it does mean the
final manuscript should avoid claiming that the mobile TFLite package has already been validated
as the same detector pipeline unless a separate model-export and label-alignment check is added.

\ChapterSubsection{Mobile Verification and Current Limitations}

The mobile application has substantial implemented functionality, but the manuscript evidence is
not yet complete. Source inspection confirms architecture, routes, repositories, persistence,
offline sync infrastructure, diagnosis calls, localization, and model-asset loading paths. It does
not by itself prove real-device performance, UI correctness on all supported screen sizes, camera
behavior on different Android devices, backend latency, or complete sync behavior across all
domains. These should be added as final manuscript evidence after screenshots, device runs, and
integration tests are collected.

\begin{table}[H]
\centering
\caption{Mobile app verification status and remaining evidence}
\label{tab:mobile-verification-limitations}
\small
\resizebox{\textwidth}{!}{%
\begin{tabular}{@{}lll@{}}
\toprule
\textbf{Area} & \textbf{Current status} & \textbf{Required final evidence} \\
\midrule
Architecture & Source-backed architecture figure generated & Keep figure and add screen-flow screenshots \\
Navigation & Route files and layouts inspected & Capture app navigation screenshots from running build \\
Scan workflow & Capture, preview, editor, analysis, result code present & Run device or emulator flow and document observed behavior \\
Offline support & SQLite, queue, sync manager implemented & Test offline-to-online recovery with real queued scans \\
Backend integration & Diagnosis client targets \texttt{/api/v1/diagnose} & Record request-response examples and latency measurements \\
TFLite deployment & Model assets and loader present & Verify exported model labels, precision behavior, and device inference time \\
Automated tests & Limited inference-engine test harness visible & Add app-level unit, integration, and workflow tests where feasible \\
\bottomrule
\end{tabular}%
}
\end{table}

Overall, the mobile app section can be written as an implemented and source-verified application
component with clear remaining evidence requirements. It should not be reduced to a placeholder:
the codebase contains a real application architecture and multiple completed user-facing
workflows. At the same time, the thesis should remain precise by distinguishing implemented source
structure from runtime validation that still needs to be captured for the final submission.

\ChapterSection{Integrated System Discussion}

\ChapterSubsection{End-to-End Integration Outcome}

The integrated Krishi Vaidya system combines mobile image capture, backend orchestration, VLM
prediction, structured advisory generation, and local farmer-facing result handling. The
implementation is not a collection of disconnected demonstrations. The mobile app contains the
scan workflow and result workflow; the backend exposes the diagnosis route; the VLM gateway
exposes readiness and inference endpoints; the advisory client transforms model context into a
schema-validated advisory; and the mobile result view consumes the returned structure for display,
history, reminders, feedback, voice output, and PDF export.

The most important integration result is the separation of responsibilities. The mobile app does
not directly call the VLM service or the Google AI service. The backend does not present raw VLM
text to the farmer. The VLM gateway does not decide how treatment should be shown in the app. Each
component has a defined boundary, and the data passed across those boundaries is validated or
normalized before it moves to the next stage. This structure reduces coupling and makes the system
easier to reason about during manuscript writing and future improvement.

\begin{table}[H]
\centering
\caption{Integrated subsystem responsibilities}
\label{tab:integrated-subsystem-responsibilities}
\small
\resizebox{\textwidth}{!}{%
\begin{tabular}{@{}lll@{}}
\toprule
\textbf{Subsystem} & \textbf{Integration responsibility} & \textbf{Evidence from implementation} \\
\midrule
Mobile app & Capture image, choose online/offline path, present result, store local state & Scan view models, diagnosis hook, API client, SQLite stores, result view model \\
Backend API & Validate upload, orchestrate VLM and advisory calls, return stable JSON & Diagnosis route, controller, validation service, diagnosis service, response types \\
VLM gateway & Validate inference request, protect inference with JWT, call vLLM, return prediction metadata & FastAPI \texttt{/ready}, \texttt{/v1/inference}, request schemas, API tests \\
VLM model & Generate crop-disease prediction text from image and prompt & Qwen2.5-VL inference gateway and validation outputs documented in VLM section \\
Advisory generator & Convert image plus VLM context into structured treatment guidance & Google AI advisory client, strict schema, grounding fallback, Zod validation \\
Local persistence & Preserve scan history, queued scans, settings, crop records, reminders, notifications & SQLite schema, repositories, offline queue, sync queue, stores \\
\bottomrule
\end{tabular}%
}
\end{table}

\ChapterSubsection{Data Contract Strengths}

The backend-to-mobile response is the strongest formal contract in the integrated system. It
contains request identifier, timestamp, disease fields, crop, confidence, confidence tier,
out-of-distribution flag, symptom descriptions, advisory arrays, processing time, and model
version. The advisory object separates treatment steps, prevention tips, safety warnings, organic
alternatives, and source URLs. This contract gives the mobile app predictable fields for display
and reduces the risk of placing unconstrained model output directly in front of the user.

The backend-to-VLM boundary is also controlled. The backend checks readiness before inference,
sends base64 image data and a deterministic prompt, limits generation length, sets temperature to
zero, disables sampling, and includes the same request identifier used by the backend response.
The VLM gateway validates authentication, image source shape, image validity, and rate limits. It
returns latency and token counts, which can later be used for deployment reporting.

The advisory-generation boundary adds another safeguard. The backend does not accept arbitrary
text as final treatment advice. The advisory client requires the Google AI response to match a
schema before returning it. It also instructs the model to produce Nepali and English fields, to
avoid fabricated citations, to include uncertainty behavior, and to avoid unsafe pesticide advice.
This does not make the output clinically or agronomically perfect, but it provides an important
engineering control around a generative step.

\ChapterSubsection{Offline and Degraded-Network Behavior}

The integrated workflow is designed for unstable connectivity. The mobile diagnosis hook has
separate behavior for online, slow, and offline network states. In a normal online state, it calls
the backend. In a slow state, it attempts the backend request but falls back after five seconds. In
an offline state, it queues the scan immediately. This design acknowledges a practical constraint:
the intended users may be in agricultural field environments where network access is unreliable.

The offline path is not only a UI message. It writes the scan into the local queue, inserts a
pending scan record into history, and creates a high-priority sync queue item. If cached advisory
data exists for the crop, the app can return a cached result while retaining the queued scan for a
future cloud diagnosis. The implementation therefore supports partial usefulness during
connectivity loss while keeping the cloud path available when the network returns.

The remaining limitation is that full two-way synchronization across all local tables is not yet
proven. The sync manager has the structure for queued operations, retry state, network monitoring,
app-resume synchronization, and periodic synchronization, but some generic table-push behavior is
still represented as local status updates until matching backend handlers are completed. This
distinction should remain clear in the final manuscript.

\ChapterSubsection{Reliability Controls Across the Integrated Path}

Reliability is handled through multiple layers rather than a single validation point. The mobile
client checks that the image URI is usable before building the request. The backend upload route
uses single-image upload middleware. The backend validation service checks file presence, file
content, MIME consistency, image format, image dimensions, and maximum size. The VLM client checks
service readiness and applies request timeouts. The VLM gateway validates JWT scope and image
payloads. The advisory client retries configured failures and validates its final JSON schema.
The mobile result view also normalizes advisory arrays to avoid downstream crashes when optional
advisory fields are missing.

\begin{table}[H]
\centering
\caption{Reliability controls in the integrated workflow}
\label{tab:integrated-reliability-controls}
\small
\resizebox{\textwidth}{!}{%
\begin{tabular}{@{}lll@{}}
\toprule
\textbf{Layer} & \textbf{Control} & \textbf{Purpose} \\
\midrule
Mobile client & Image URI validation, retry interceptor, slow/offline fallback & Avoid invalid local requests and preserve workflow under poor connectivity \\
Backend upload & Multipart single-image route and upload middleware & Restrict diagnosis route to one uploaded image input \\
Backend validation & MIME/header match, file size, image dimensions, crop-name validation & Reject corrupt, unsupported, or unsafe request data before model calls \\
Backend VLM client & Readiness probe, timeout, bearer or RS256 JWT authentication & Protect inference service and avoid uncontrolled hanging requests \\
VLM gateway & JWT scope check, image validation, rate limiting, vLLM health check & Protect model endpoint and reject invalid inference inputs \\
Advisory client & Retry, optional grounding fallback, strict schema validation & Prevent invalid generated advisory JSON from reaching mobile clients \\
Mobile result layer & Safe advisory defaults, local history save, queued state handling & Keep result screens stable across cloud, cached, and queued outputs \\
\bottomrule
\end{tabular}%
}
\end{table}

These controls are valuable because the system depends on machine-learning and generative-model
components. A model can return malformed text, an external service can be unavailable, or a mobile
network can fail mid-request. The integration design does not remove those risks, but it limits
how far an invalid or unavailable component can propagate through the user-facing workflow.

\ChapterSubsection{Integration Limitations and Evidence Gaps}

The current integration evidence is source-backed and architecture-backed, but it is not yet a
complete deployment evaluation. The manuscript now has the implemented route contracts, generated
integration diagram, method description, and integration discussion. However, the final thesis
should still include live evidence from the complete running system. The most important missing
items are a real mobile-to-backend diagnosis request, captured response JSON, backend logs or
screenshots, VLM gateway readiness and inference logs, mobile result screenshots, offline queue
demonstration, and measured latency from image submission to final result.

There are also two technical alignment issues that must remain visible. First, VLM evaluation in
this manuscript shows a high parse-error rate and low prediction accuracy on the independent test
set. The backend advisory layer reduces the risk of exposing raw unstructured VLM output, but it
does not by itself prove that the current VLM is diagnostically reliable. Second, the mobile
TFLite asset labels visible in the source differ from the YOLO manuscript detector taxonomy. This
requires a final export-label alignment check before claiming that the packaged mobile model is
the same detector described in the YOLO section.

\begin{table}[H]
\centering
\caption{Integrated system evidence status}
\label{tab:integrated-evidence-status}
\small
\resizebox{\textwidth}{!}{%
\begin{tabular}{@{}lll@{}}
\toprule
\textbf{Evidence item} & \textbf{Current status} & \textbf{Required before final claim} \\
\midrule
Architecture and control flow & Generated and documented from source & Review figure placement and wording in final PDF \\
Mobile-to-backend route contract & Verified from mobile API client and backend route/controller & Capture real request and response from running services \\
Backend-to-VLM contract & Verified from backend VLM client and FastAPI gateway & Capture readiness and inference logs from running VLM gateway \\
Advisory schema contract & Verified from advisory client and schema validation & Capture at least one valid generated advisory response \\
Offline queue behavior & Verified from queue manager and sync infrastructure & Demonstrate offline queue and later recovery in app runtime \\
Runtime performance & Not yet measured end to end & Measure mobile request latency, backend processing time, and VLM latency \\
Model alignment & Partially documented with caveats & Reconcile mobile TFLite labels with YOLO training/export assets \\
\bottomrule
\end{tabular}%
}
\end{table}

\ChapterSubsection{Overall Integrated System Interpretation}

The integrated system should be described as a substantial implemented prototype with a complete
source-backed diagnosis pathway and clear remaining runtime-evidence requirements. Its strongest
engineering qualities are modular boundaries, explicit validation, structured response contracts,
offline-aware mobile behavior, and farmer-facing result handling. Its main weaknesses are not
architectural absence but incomplete final verification: live end-to-end screenshots, latency
measurements, mobile model alignment, and broader automated integration tests still need to be
added.

This interpretation is important for manuscript accuracy. Overclaiming production readiness would
be unsupported, but under-describing the integration would also be inaccurate. The source code
shows that Krishi Vaidya is not merely a YOLO model, a VLM experiment, a backend, or a mobile app
in isolation. It is an integrated agricultural advisory system whose final thesis narrative should
present both the completed engineering structure and the remaining verification work with equal
precision.

\ChapterSection{Cross-System Results Discussion and Limitations}

\ChapterSubsection{Synthesis of Major Results}

The main result of Krishi Vaidya is an implemented multi-component agricultural advisory system
with substantial evidence at the model, backend, mobile, and integration layers. The YOLO
subsystem provides the strongest quantitative model result. It has training curves, threshold
analysis, deployment comparisons, per-class metrics, confusion-matrix outputs, and error-analysis
artifacts. The VLM subsystem provides a complete fine-tuning and evaluation pipeline, but its
prediction results reveal a serious reliability limitation. The backend subsystem provides a
structured API layer with domain models, upload validation, authentication/session controls,
diagnosis orchestration, advisory generation, and scan-domain tests. The mobile app provides a
farmer-facing interface with scan workflow, offline-aware diagnosis behavior, crop management,
history, settings, localization, and result actions. The integration workflow connects these
components through a source-verified diagnosis path.

The strongest overall contribution is therefore not a single model score. It is the construction
of a full crop-health support pipeline that begins with field image capture and ends with a
structured advisory report. This matters because many prototype systems stop at a classifier
output. Krishi Vaidya contains the additional engineering needed to make the output usable:
backend validation, structured response assembly, mobile presentation, scan history, reminders,
offline queueing, and farmer-facing language support.

\begin{table}[H]
\centering
\caption{Cross-system result synthesis}
\label{tab:cross-system-result-synthesis}
\small
\resizebox{\textwidth}{!}{%
\begin{tabular}{@{}lll@{}}
\toprule
\textbf{System area} & \textbf{Strongest result} & \textbf{Main caution} \\
\midrule
YOLO detector & Quantitative detection and deployment evidence is complete enough for manuscript use & Mobile package alignment and physical-device profiling still need final verification \\
VLM model & Dataset export, fine-tuning, and validation pipeline are complete and documented & Prediction accuracy and parseability are currently weak on the independent test set \\
Backend API & Structured diagnosis orchestration, domain models, validation, and scan tests exist & Broader route and diagnosis-integration tests remain incomplete \\
Mobile app & Complete source-backed farmer workflow and offline-aware structure exist & Final screenshots, device tests, and runtime measurements are missing \\
Integration & End-to-end source-verified diagnosis path exists across mobile, backend, VLM, and advisory layers & Live request-response and latency evidence still needs to be captured \\
\bottomrule
\end{tabular}%
}
\end{table}

\ChapterSubsection{Model-Level Interpretation}

The YOLO model results indicate that the object-detection component is the most mature machine
learning part of the project. Its metrics are not perfect, but they are supported by multiple
forms of evidence: training dynamics, robust validation, threshold analysis, quantized deployment
comparison, latency measurements, and error analysis. This allows the manuscript to discuss
realistic deployment tradeoffs instead of only reporting a single final accuracy number.

The VLM results must be interpreted differently. The VLM work is substantial because it includes
dataset construction, Qwen-format export, LoRA fine-tuning, quantized training, inference
infrastructure, dataset validation, and prediction evaluation. However, the independent prediction
metrics show that the current VLM is not yet reliable enough to be presented as a final
diagnostic authority. The high parse-error count is especially important because the backend and
mobile app depend on structured fields. A model that often fails to produce parseable output can
still be useful as a research artifact, but it must not be overclaimed as a dependable diagnosis
engine.

The correct interpretation is that the project has one relatively stronger model component and
one exploratory but incomplete VLM component. The YOLO section supports confidence in the
object-detection work. The VLM section supports confidence in the pipeline and exposes the next
research problem: improving structured multimodal disease prediction.

\ChapterSubsection{System-Level Interpretation}

The backend and mobile results show that the project was engineered as a system rather than as an
isolated notebook experiment. The backend provides request validation, route organization,
repository-backed domain logic, image upload support, token/session handling, scan lifecycle
management, crop/task/note/advisory records, VLM inference bridging, and advisory-schema
validation. The mobile app provides the user-facing workflow required to apply those backend
features in a field setting.

The important system-level finding is that the engineering structure is ahead of the final runtime
evidence. The source code shows a serious architecture, but the final manuscript still needs
runtime captures to prove how the system behaves when all services are running together. This is a
normal distinction in software-engineering research: implemented structure and runtime validation
are related, but they are not the same evidence. The manuscript should use both, and should never
replace missing runtime evidence with assumptions.

\ChapterSubsection{Reliability and Safety Interpretation}

Krishi Vaidya includes several reliability controls that are appropriate for an agricultural
advisory system. Images are validated before model calls. Backend diagnosis uses request IDs and
processing-time measurement. The VLM service is checked for readiness before inference. The VLM
gateway uses JWT scope validation. The advisory generator is required to return schema-valid JSON.
The mobile app handles cloud, cached, and queued outcomes. The result view uses structured fields
instead of raw model text.

These controls improve engineering safety, but they do not solve every agronomic safety problem.
Model uncertainty, class imbalance, field-image variability, pesticide recommendation risk,
out-of-distribution images, and incorrect crop hints can still affect final advice. The system
therefore should present confidence, uncertainty, retake guidance, safety warnings, and expert
consultation guidance where appropriate. The backend already includes confidence-tier and
out-of-distribution fields, which should remain central in the final app presentation and
manuscript claims.

\ChapterSubsection{Current Limitations}

The current limitations fall into four groups: model limitations, validation limitations,
integration limitations, and manuscript-evidence limitations.

Model limitations are most visible in the VLM results. The VLM has low test accuracy and a high
number of parse errors. This limits the strength of claims that can be made about automatic
disease diagnosis. The YOLO model is stronger, but its mobile deployment path still requires final
alignment between manuscript detector taxonomy and packaged mobile model labels.

Validation limitations are visible in the backend and mobile app. Backend scan tests are useful,
but equivalent tests are not yet available for every route group and external diagnosis branch.
The mobile app has source-backed architecture and an inference harness, but not a complete
recorded emulator or physical-device validation suite. Integration validation is source-verified
but not yet demonstrated through a captured live end-to-end request.

Integration limitations are tied to external services and runtime configuration. The backend
diagnosis route depends on the VLM service and Google AI advisory service being reachable and
properly configured. The advisory generator must return schema-valid JSON. The mobile app must be
configured with the correct API base URL and must handle network failure correctly. Offline sync
infrastructure exists, but complete two-way synchronization for every local table is not yet
proven.

Manuscript-evidence limitations are practical but important. The final report still needs
screenshots, live request-response examples, latency measurements, deployment notes, and runtime
logs. Without those artifacts, the manuscript can accurately claim implemented architecture and
source-verified workflows, but it should not claim full production deployment or field-tested
performance.

\begin{table}[H]
\centering
\caption{Cross-system limitations and required future evidence}
\label{tab:cross-system-limitations}
\small
\resizebox{\textwidth}{!}{%
\begin{tabular}{@{}lll@{}}
\toprule
\textbf{Limitation group} & \textbf{Current limitation} & \textbf{Required action} \\
\midrule
VLM reliability & Low prediction accuracy and high parse-error count & Improve prompt format, constrained decoding, dataset balance, and evaluation feedback loop \\
YOLO deployment alignment & Mobile TFLite label list differs from manuscript detector taxonomy & Verify export source, label map, and mobile model package before final claim \\
Backend test coverage & Scan domain is tested more deeply than other route domains & Add auth, crop, task, note, advisory, diagnosis, and upload-validation tests \\
Mobile validation & No complete recorded emulator/physical-device workflow evidence yet & Capture screenshots, scan flow, result flow, offline queue, and TFLite behavior \\
Integration evidence & Source-verified but not yet live end-to-end demonstrated in manuscript & Capture real request-response logs and latency across mobile, backend, VLM, and advisory layers \\
Offline synchronization & Queue and sync infrastructure exists, but full table sync is incomplete & Add backend handlers and demonstrate offline-to-online recovery \\
Agronomic safety & Generated advice can still be uncertain or unsafe if not constrained and reviewed & Preserve confidence tiers, safety warnings, source links, and expert consultation guidance \\
\bottomrule
\end{tabular}%
}
\end{table}

\ChapterSubsection{Implications for Final Manuscript Claims}

The final manuscript should make strong claims where the evidence is strong and cautious claims
where the evidence is incomplete. It can strongly claim that the project implements a modular
Krishi Vaidya system with YOLO, VLM, backend, mobile, and integration components. It can strongly
claim that YOLO and VLM datasets and outputs have been analyzed and documented using available
figures and tables. It can strongly claim that the backend and mobile application contain
substantial implemented architecture.

It should not claim that the current VLM model is already highly accurate. It should not claim
that the mobile app has been field-tested unless those tests are actually performed. It should not
claim full production readiness until deployment, latency, screenshots, runtime logs, and broader
automated tests are added. This disciplined boundary will make the thesis stronger, not weaker,
because it shows that the project team understands the difference between implementation,
evaluation, and deployment.
